\section{Endpoint polynomial proofs, certificates and covering scope}
\label{sec:endpoint-details}

The two possible low integer roots have different arithmetic constraints.
The endpoint two is confined to a linear region in the minimum exponent;
this gives a new nonintegral tail for even exponent triples.
The shifted odd-root distributions extend the same tail to two mixed-parity
patterns below.

\begin{proposition}
\label{thm:endpoint-two-bound}
Let $4\leq a\leq b\leq c$, and put
\[
R(a)=\frac{7a^2-2a+8}{a+2}=7a-16+\frac{40}{a+2}.
\]
If $h_C(2)=0$, then $c<R(a)$. In fact, $c\geq R(a)$ implies
$h_C(2)>0$.
\end{proposition}

\begin{proof}
Fix $a,c$ and expand the symmetric endpoint polynomial in $b$:
\[
h_C(2)=A(a,c)b^3+B(a,c)b^2+D(a,c)b+E(a,c),
\]
where
\begin{align*}
A(a,c)&=ac(a+c-1)+2a(a-1)+2c(c-1)-4,\\
B(a,c)&=c^2\bigl((a+2)c-(7a^2-2a+8)\bigr)\\
&\quad +(a^3+2a^2-6a-4)c+2(a-2)(a^2-2a-6),\\
D(a,c)&=(a-2)(c-2)\bigl(a^2c+a^2+ac^2+6ac+4a+c^2+4c-12\bigr),\\
E(a,c)&=2(a-2)(c-2)(a+c-2)(ac+a+c-2).
\end{align*}
For $a,c\geq4$, the coefficients $A,D,E$ are strictly positive.
Both $a^3+2a^2-6a-4$ and $a^2-2a-6$ are positive at $a=4$
and increasing on $[4,\infty)$. Thus $c\geq R(a)$ also gives $B>0$,
which proves the strict endpoint sign. Conversely, an endpoint zero
requires $B<0$, so $(a+2)c<7a^2-2a+8$.
\end{proof}

\begin{proposition}
\label{cor:even-linear-tail}
If $4\leq a\leq b\leq c$ are all even and $c\geq R(a)$,
then $G_{p^a q^b r^c}$ is not Laplacian integral.
For minimum exponent at least eight, $c\geq7a-12$ suffices.
\end{proposition}

\begin{proof}
Every entry of $C$ is even, so every integer eigenvalue of $C$ is even
by the rational-root theorem applied to $C/2$. In particular, one is
not an eigenvalue. Proposition~\ref{thm:endpoint-two-bound} excludes two,
and Theorem~\ref{thm:uniform-low-root} provides a noninteger root in
$(0,3)$. Equivalently, $h_C(0)<0<h_C(2)$ directly yields a root in
$(0,2)$, which cannot equal one. Its graph lift is noninteger.
For $a\geq8$, $40/(a+2)\leq4$ gives $R(a)\leq7a-12$.
\end{proof}

The triples $(8,10,44)$, $(10,14,58)$ and $(40,42,266)$ meet this
criterion with gcd two. Their largest exponents are below the earlier
quadratic cutoff and their spans are outside the earlier balanced region.
The already settled repeated triple $(10,10,12)$ has $h_C(2)=0$ and
obeys the strict bound; an integer endpoint does not establish integrality.

\begin{proposition}
\label{cor:mixed-linear-tail}
Let $4\leq m\leq d\leq M$ be the ordered exponents of a triple that is
a permutation of $(1,3,2)$ or $(3,3,0)$ modulo four.
If $M\geq R(m)$, the ideal intersection graph is nonintegral.
For $m\geq8$, the simpler bound $M\geq7m-12$ suffices.
\end{proposition}

\begin{proof}
Under the integer-spectrum assumption,
Lemma~\ref{lem:mixed-endpoint-one} excludes one as a quotient root.
Theorem~\ref{thm:uniform-low-root} then forces two to be a root.
But Proposition~\ref{thm:endpoint-two-bound} gives $h_C(2)>0$ for
$M\geq R(m)$, a contradiction. The simpler bound follows from
$R(m)\leq7m-12$ for $m\geq8$.
\end{proof}

Thus the linear upper bound also restricts these two mixed-parity patterns.
For example, the residue-role triples $(25,27,758)$ and $(29,27,882)$
satisfy the arithmetic conditions of Proposition~\ref{thm:clustered-odd-roots},
but their largest exponents exceed this linear bound. The root-distribution
argument therefore settles these compatible controls without a further
congruence lift. The triple $(19,27,140)$ illustrates the other pattern.
The variables $m,d,M$ identify size order, independently of residue roles.

The size bound confines one surface geometrically. The following
constant-term conditions give finite candidate sets for either surface
when two exponents are fixed.

\begin{proposition}
\label{prop:endpoint-divisors}
For fixed $a,b\geq4$, any positive integer $c$ on an endpoint-zero
surface satisfies
\begin{align*}
h_C(1)=0&\ \Longrightarrow\ c\mid(a-1)(b-1)(a+b-1)(ab+a+b-1),\\
h_C(2)=0&\ \Longrightarrow\ c\mid2(a-2)(b-2)(a+b-2)(ab+a+b-2).
\end{align*}
\end{proposition}

\begin{proof}
Both endpoints are integer cubics in $c$. Their constant terms are
the displayed nonzero products. Reduce each endpoint equation modulo $c$.
\end{proof}

For a fixed pair this gives finite divisor candidates, which must still
satisfy the endpoint equation. Reza's requested diagnostic for middle
exponents through fifteen now has a complete independent certificate.

\begin{proposition}
\label{cor:middle-fifteen}
Every positive ordered three-prime exponent triple $a\leq b\leq c$
with $b\leq15$ gives a nonintegral ideal intersection graph.
This conclusion uses the complete finite endpoint certificate below.
\end{proposition}

\begin{proof}
Repeated exponents are covered by Theorem~\ref{thm:repeated}, and
minimum exponents one, two and three by Theorems~\ref{thm:unit},
\ref{thm:minimum-two} and~\ref{thm:minimum-three}, respectively.
It remains to consider $4\leq a<b\leq15$ and $c>b$.

For each of these $66$ pairs, both endpoint polynomials are integer
cubics in $c$ with the positive constant terms in
Proposition~\ref{prop:endpoint-divisors}. The
\href{https://github.com/the-omega-institute/ideal-intersection-laplacian/blob/10ff01f0625c7264b8bc5bb33b1a9d7dcb88a93c/results/endpoint-surfaces-verification.json}{complete endpoint certificate}
records all $132$
coefficient lists and exhausts their $8\,527$ positive divisors using
integer trial division. Independent integer Horner evaluation at every
one of the $7\,585$ divisors above $b$ gives no zero. Thus neither
endpoint vanishes for any integer $c>b$, with no imposed upper bound
on $c$. Theorem~\ref{thm:uniform-low-root} supplies a positive root
in $(0,3)$; it can equal neither one nor two, proving nonintegrality.
\end{proof}

The reflected quotient in this diagnostic uses
$N=|V(H)|=a+b+c+ab+ac+bc$, not the full graph vertex count.
The inclusive requested bound $b\leq15$ and exact rational linear-factor
extraction are cross-checked against the independent divisor certificate.
The empty lists settle this pair range; they do not prove either
endpoint surface globally empty.

\begin{proposition}
\label{thm:modulo-three-endpoints}
Every positive three-prime exponent triple congruent to $(2,2,2)$,
or to a permutation of $(1,2,2)$, modulo three gives a nonintegral
ideal intersection graph.
\end{proposition}

\begin{proof}
For these two residue patterns, reduction of the general complement
quintic gives, respectively,
\[
h_C(x)\equiv x^5,\qquad
h_C(x)\equiv(x^2+1)(x^3+x^2-x+1)\pmod3.
\]
The corresponding pairs $(h_C(1),h_C(2))$ are $(1,2)$ and $(1,1)$
modulo three, so neither integer endpoint is a root. Minimum exponents
at most three are already covered. Otherwise
Theorem~\ref{thm:uniform-low-root} supplies a positive root in $(0,3)$,
which cannot be an integer. Its graph lift is noninteger as well.
For the second pattern, the irreducible factor $x^2+1$ also directly
rules out a completely integral quotient spectrum.
\end{proof}

For example, $(8,17,22)$ and $(8,17,23)$ satisfy the new conditions
with gcd one and unequal $2$-adic valuations. Complete endpoint tables
modulo two and three, checked against independent integer determinants,
are given in the \href{https://github.com/the-omega-institute/ideal-intersection-laplacian/blob/2df3cffd062ce0b28fc2040a6203fd4e7282a725/results/endpoint-residues.json}{residue certificate}.
Every pair $(a,b)$ modulo either prime has at least one $c$ root on
each surface; retaining the full root sets, rather than only root-free
pairs, is necessary to detect these exclusions.

\begin{remark}
Endpoint tests at a fixed finite set of moduli cannot by themselves
prove either surface empty on all fully distinct positive triples.
Let $M$ be their least common multiple. Since $(9,9,136)$ has
$h_C(1)=0$, the strictly ordered triples
$(9+kM,9+(k+1)M,136+(k+2)M)$ pass every selected endpoint-one test.
The known endpoint-two zero $(10,10,12)$ similarly gives survivors
$(10+kM,10+(k+1)M,12+(k+2)M)$. For large $k$ their middle exponents
exceed fifteen. These are congruence survivors, not asserted integer
zeros. These fixed-span tails are excluded by
Corollary~\ref{cor:balanced-span}. To retain the geometric restrictions,
take $L$ divisible by $M$ and eight and any $d\geq32$ divisible by $L$.
The triples $(9+d,9+2d,136+4d)$ and $(10+d,10+2d,12+4d)$ pass the
respective endpoint tests, have gcd one and two, and satisfy both upper
bounds on the maximum while failing the sufficient inertia inequality.
The substitution identities are given in the accompanying
\href{https://github.com/the-omega-institute/ideal-intersection-laplacian/blob/nikandish-math-patch-1/notes/endpoint-congruence-scope.md}{scope note}.
This concerns endpoint-value tests alone; another-root obstruction or
further arithmetic information can still exclude these candidates.
\end{remark}

Theorem~\ref{thm:endpoint-two-completion} completes the endpoint-two
alternative, including the all-even triples formerly left inside this
linear region. The remaining mixed-parity surface $h_C(1)=0$ still
needs another obstruction.
